\clearpage
\chapter*{Erratum à l'exposé n\textsuperscript{o} 232}
\addcontentsline{toc}{chapter}{Erratum à l'exposé n\textsuperscript{o} 232}
\src{303}

\noindent\textbf{7. TDTE V : Les schémas de Picard. Théorèmes d'existence.}

\begin{center}
[Séminaire Bourbaki, t. 14, 1961/62, n\textsuperscript{o} 232, 19 p.]
\end{center}

\noindent\textbf{Page 7, remarque 3.3.} --- La question soulevée ici a été
résolue affirmativement par MUMFORD (voir commentaires à l'exposé suivant).

\noindent\textbf{Page 13, remarque 5.2.} --- Comme je le signale au début de
l'exposé suivant, la conjecture d'existence avancée ici est fausse, cependant
MUMFORD arrive à prouver par ses méthodes un théorème légèrement plus faible.

\noindent\textbf{Page 14, ligne 7.} --- Après « complet », ajouter « à corps
résiduel algébriquement clos ». Le contre-exemple de MUMFORD qu'on vient de
signaler montre d'ailleurs que cette restriction est indispensable.

\noindent\textbf{Page 17, remarque 6.6.} --- Comme nous le signalons au début
de l'exposé suivant, la question soulevée ici vient d'être résolue par
l'affirmative par MURRE.

\noindent\textbf{Page 18, remarque 6.7.} --- La question soulevée ici se résout
par l'affirmative (cf. dernier alinéa des commentaires à l'exposé suivant).

\noindent\textbf{Page 18, remarque 6.8, 4\textsuperscript{e} ligne.} --- Au lieu
de « régulier », lire « simple sur $k$ ».

\par\smallskip
\begingroup\footnotesize
\endgroup
